\newcounter{RDW} \setcounter{RDW}{1}
\newcounter{RDR} \setcounter{RDR}{1}
\newcounter{RDE} \setcounter{RDE}{1}
\newcounter{RDS} \setcounter{RDS}{1}
\newcounter{RF} \setcounter{RF}{1}

% ======================= %
%     ENVIAR MENSAJE      %
% ======================= %
\section{Enviar mensaje}
\textbf{\textcolor{rojo}{\underline{RF5.\theRF:}}}. Enviar mensaje: Permite que un usuario envíe un mensaje de un usuario a otro. Los datos necesarios son el usuario 
emisor y el receptor (Cadenas de caracteres de 40 caracteres), el mensaje (cadena de caracteres de 100 máximo) y a parte, la base de datos almacenará un bit de 
visualización (booleano) por defecto en 0, la hora y fecha de envío y el ID del mensaje (cadena de 50 caracteres). Se requiere que tanto el emisor como el 
receptor existan, además de que no sean el mismo usuario. Por otro lado, también se requiere que ambos sean amigos y que el mensaje que va a ser enviado no 
puede estar vacío.\\[1pt]

\textbf{\textcolor{rojo}{Entrada:}} Agente externo: usuario. Acción: solicitar envío de un mensaje. 
Requisito de datos de entrada \textcolor{azul}{RDE5.\theRDE}. \\[6pt]

\textbf{\textcolor{rojo}{BD:}} Requisito de datos de escritura \textcolor{azul}{RDW5.\theRDW} y de lectura \textcolor{azul}{RDR5.\theRDR}. \\[6pt]

\textbf{\textcolor{rojo}{Salida:}} Agente externo: usuario. Acción: confirmación resultado. 
Requisito de datos de salida: ninguno. \\[12pt]

\textcolor{azul}{\textbf{RDE5.\theRDE:}} Mensaje y usuario al que va dirigido \\
\hspace*{1cm} ID Usuario: Cadena de caracteres (40) \\
\hspace*{1cm} Mensaje: Cadena de caracteres (100) \\[12pt]

\textcolor{azul}{\textbf{RDR5.\theRDR:}} Comprobación de que existen ambos usuarios y si son amigos entre sí (Que un usuario sea amigo de otro no implica que sea al revés). \\
\hspace*{1cm} ID Usuario Emisor: Cadena de caracteres (40) \\
\hspace*{1cm} ID Usuario Receptor: Cadena de caracteres (40) \\[12pt]

\textcolor{azul}{\textbf{RDW5.\theRDW:}} Datos del envío de mensaje. \\
\hspace*{1cm} ID Usuario Emisor: Cadena de caracteres (40) \\
\hspace*{1cm} ID Usuario Receptor: Cadena de caracteres (40) \\
\hspace*{1cm} Fecha y hora del envío: Variable tipo fecha \\
\hspace*{1cm} Bit de visto: Booleano (1)\\
\hspace*{1cm} ID mensaje: Cadena de caracteres (50) \\
\hspace*{1cm} Mensaje: Cadena de caracteres (100) \\[12pt]


\newcounter{RS} \setcounter{RS}{1}
\textbf{\textcolor{verde}{Restricciones Semánticas:}} \\[6pt]
\textbf{\textcolor{verde}{\underline{RS5.\theRF.\theRS:}}} Los mensajes sólo pueden enviarse entre usuarios que son amigos. \\
\textbf{RF:} RF5.\theRF. \textbf{RD(s):} RDE5.\theRDE. RDR5.\theRDR \\
\hspace*{1cm} \textit{Descripción:} Si se desea enviar un mensaje a un usuario del que no se es amigo, se devolverá un error indicando que se necesita ser amigos para poder enviar un mensaje y no se envía el mensaje. \\[6pt]
\stepcounter{RS}

\textbf{\textcolor{verde}{\underline{RS5.\theRF.\theRS:}}} Los mensajes no se pueden enviar al propio usuario. \\
\textbf{RF:} RF5.\theRF. \textbf{RD(s):} RDE5.\theRDE. \\
\hspace*{1cm} \textit{Descripción:} Si se desea enviar un mensaje al propio usuario, se devolverá un error indicando que no se puede enviar un mensaje a uno mismo y no se envía el mensaje. \\[6pt]
\stepcounter{RS}

\textbf{\textcolor{verde}{\underline{RS5.\theRF.\theRS:}}} Para el envío de mensajes, el mensaje debe tener al menos un carácter. \\
\textbf{RF:} RF5.\theRF. \textbf{RD(s):} RDE5.\theRDE. \\
\hspace*{1cm} \textit{Descripción:} Si se desea enviar un mensaje, dicho mensaje no puede estar vacío, de no ser así, se muestra un error y no se envía el mensaje. \\ [6pt]
\stepcounter{RS}

\textbf{\textcolor{verde}{\underline{RS5.\theRF.\theRS:}}} Para el envío de mensajes, ambos usuarios deben existir. \\
\textbf{RF:} RF5.\theRF. \textbf{RD(s):} RDE5.\theRDE. RDR5.\theRDR \\
\hspace*{1cm} \textit{Descripción:} Si se desea enviar un mensaje, los usuarios deben de existir, de no ser así damos un mensaje de error e indicamos que el usuario no existe. \\ [6pt]
\stepcounter{RS}

\stepcounter{RF}
\stepcounter{RDE}
\stepcounter{RDW}
\stepcounter{RDR}

% ======================= %
%    ELIMINAR MENSAJE     %
% ======================= %
\section{Eliminar mensaje}
\textbf{\textcolor{rojo}{\underline{RF5.\theRF:}}}. Eliminar mensaje: Permite que un usuario elimine un mensaje de texto que ha enviado, de forma que ni dicho usuario ni el receptor pueden volver a acceder al mensaje. Se requiere del usuario emisor y del receptor (cadena de 40 caracteres) además del ID del mensaje por borrar (cadena de 50 caracteres). Para ello, se requiere de que el mensaje exista\\[6pt]

\textbf{\textcolor{rojo}{Entrada:}} Agente externo: usuario. Acción: solicitar la eliminación de un mensaje. 
Requisito de datos de entrada: \textcolor{azul}{RDE5.\theRDE.}\\[6pt]

\textbf{\textcolor{rojo}{BD:}} Requisito de datos de escritura \textcolor{azul}{RDW5.\theRDW} y de lectura \textcolor{azul}{RDR5.\theRDR}. \\[6pt]

\textbf{\textcolor{rojo}{Salida:}} Agente externo: usuario. Acción: confirmación resultado. 
Requisito de datos de salida: ninguno. \\[12pt]

\textcolor{azul}{\textbf{RDE5.\theRDE:}} Datos del mensaje a borrar\\
\hspace*{1cm} ID Usuario Emisor: Cadena de caracteres (40) \\
\hspace*{1cm} ID Usuario Receptor: Cadena de caracteres(40) \\
\hspace*{1cm} ID Mensaje: Cadena de caracteres (50) \\[12pt]

\textcolor{azul}{\textbf{RDR5.\theRDR:}} Comprobación de que existe el mensaje a borrar. \\
\hspace*{1cm} ID Usuario Emisor: Cadena de caracteres (40) \\
\hspace*{1cm} ID Usuario Receptor: Cadena de caracteres (40) \\
\hspace*{1cm} ID Mensaje: Cadena de caracteres (50) \\[12pt]

\textcolor{azul}{\textbf{RDW5.\theRDW:}} Eliminación del mensaje con los mismos datos que \textcolor{azul}{\textbf{RDE5.\theRDE}} además de:
\hspace*{1cm} Mensaje: Cadena de caracteres (50) \\
\hspace*{1cm} Bit de visualización: Booleano (1) \\
\hspace*{1cm} Fecha de envío: Variable tipo fecha \\
\hspace*{1cm} Hora de envío: Variable tipo hora. \\[12pt]

\setcounter{RS}{1}
\textbf{\textcolor{verde}{Restricciones Semánticas:}} \\[6pt]
\textbf{\textcolor{verde}{\underline{RS5.\theRF.\theRS:}}} Para eliminar un mensaje, este debe de existir. \\
\textbf{RF:} RF5.\theRF. \textbf{RD(s):} RDE5.\theRDE, RDR5.\theRDR \\
\hspace*{1cm} \textit{Descripción:} Si se desea eliminar un mensaje, este debe de existir, de no ser así, se mostrará un mensaje de error indicando que el mensaje no existe y no se modificará la lista de mensajes. \\ [6pt]
\stepcounter{RS}

\stepcounter{RF}
\stepcounter{RDE}
\stepcounter{RDW}
\stepcounter{RDR}
% ======================= %
%  LISTAR CONVERSACIONES  %
% ======================= %
\section{Listar usuarios}
\textbf{\textcolor{rojo}{\underline{RF5.\theRF:}}}. Listar usuarios: Permite a un usuario ver los usuarios con los que puede tener mensajes (con los que es amigo recíprocamente) y no están borrados, ordenados temporalmente 
y por visualización (Primero aquellos con mensajes no visualizados después visualizados, cada uno ordenado temporalmente de más a menos reciente). Se puede solicitar ver los usuarios archivados como los que no, para poder 
gestionar ambos grupos.\\[6pt]

\textbf{\textcolor{rojo}{Entrada:}} Agente externo: usuario. Acción: solicitar la visualización de los usuarios con los que puede conversar. 
Requisito de datos de entrada \textcolor{azul}{RDE5.\theRDE}. \\[6pt]

\textbf{\textcolor{rojo}{BD:}} Requisito de datos de lectura \textcolor{azul}{RDR5.\theRDR}. \\[6pt]

\textbf{\textcolor{rojo}{Salida:}} Agente externo: usuario. Acción: mostrar los usuarios con los que tiene mensajes. 
Requisito de datos de salida: \textcolor{azul}{RDS5.\theRDS} \\[12pt]

\textcolor{azul}{\textbf{RDE5.\theRDE:}} Grupo de usuarios que se quiere mostrar archivados/no archivados\\
\hspace*{1cm} Bit elección (indica si se quiere ver el grupo de conversaciones archivadas o las no archivadas True/False): Booleano (1) \\[12pt]

\textcolor{azul}{\textbf{RDR5.\theRDR:}} Lista de usuarios archivados o no archivados, con los que puede tener una conversación, ordenados según lo especificado, con el formato:\\
\hspace*{1cm} ID Usuario: Cadena de caracteres (40) \\
\hspace*{1cm} Bit elección (indica si se quiere ver el grupo de conversaciones archivadas o las no archivadas True/False): Booleano (1) \\[12pt]

\textcolor{azul}{\textbf{RDS5.\theRDS:}} Lista de usuarios archivados o no con los que puede tener una conversación. \\
\hspace*{1cm} Nombre Usuario con ID usuario: Cadena de caracteres (20) [12pt]

\setcounter{RS}{1}
\textbf{\textcolor{verde}{Restricciones Semánticas:}} \\[6pt]
\textbf{\textcolor{verde}{\underline{RS5.\theRF.\theRS:}}} Sólo se muestran los usuarios archivados o no archivados. \\
\textbf{RF:} RF5.\theRF. \textbf{RD(s):} RDR5.\theRDR \\
\hspace*{1cm} \textit{Descripción:} Los usaurios de la opción no escogida no se muestran, de no ser así, se muestra un error indicandolo. \\ [6pt]
\stepcounter{RS}

\stepcounter{RF}
\stepcounter{RDR}
\stepcounter{RDS}
\stepcounter{RDE}
% ======================= %
%     VISUALIZAR CHAT     %
% ======================= %
\section{Visualizar mensajes}
\textbf{\textcolor{rojo}{\underline{RF5.\theRF:}}}. Visualizar mensajes: Permite que a usuario ver los mensajes enviados y recibidos entre un usuario y otro, 
ordenados de forma temporal, de más a menos reciente. Se requiere del id del usuario del que se desea ver la conversación, además del propio usuario (Cadenas 
de 40 caracteres). Actualizamos el bit de visualización a "Visto" tras esta acción en la conversación correspondiente.\\[6pt]

\textbf{\textcolor{rojo}{Entrada:}} Agente externo: usuario. Acción: solicitar la visualización de un chat(mensajes entre un usuario y otro). 
Requisito de datos de entrada \textcolor{azul}{RDE5.\theRDE}. \\[6pt]

\textbf{\textcolor{rojo}{BD:}} Requisito de datos de lectura \textcolor{azul}{RDR5.\theRDR} y de escritura \textcolor{azul}{RDW5.\theRDW}. \\[6pt]

\textbf{\textcolor{rojo}{Salida:}} Agente externo: usuario. Acción: confirmación resultado. 
Requisito de datos de salida: \textcolor{azul}{RDS5.\theRDS} \\[12pt]

\textcolor{azul}{\textbf{RDE5.\theRDE:}} Datos del usuario del que se desea ver los mensajes entre ambos\\
\hspace*{1cm} ID Usuario solicitado: Cadena de caracteres (40) \\[12pt]

\textcolor{azul}{\textbf{RDR5.\theRDR:}} Lista de mensajes entre ambos usuarios\\
\hspace*{1cm} ID Usuario solicitante: Cadena de caracteres (40) \\
\hspace*{1cm} ID Usuario solicitado: Cadena de caracteres (40) \\[12pt]

\textcolor{azul}{\textbf{RDW5.\theRDW:}} Actualizamos los mensajes como leidos por el usuario poniendolos a 1:\\
\hspace*{1cm} ID Usuario lector: Cadena de caracteres (40) \\
\hspace*{1cm} ID del otro usuario: Cadena de caracteres (40) \\
\hspace*{1cm} Bit de visto: Booleano (1) \\[12pt]

\textcolor{azul}{\textbf{RDS5.\theRDS:}} Lista de mensajes entre ambos usuarios, con formato:\\
\hspace*{1cm} ID Usuario emisor: Cadena de caracteres (40) \\
\hspace*{1cm} Mensaje: Cadena de caracteres (50) \\[12pt]

\setcounter{RS}{1}
\textbf{\textcolor{verde}{Restricciones Semánticas:}} \\[6pt]
\textbf{\textcolor{verde}{\underline{RS5.\theRF.\theRS:}}} Para visualizar los mensajes debe de existir el usuario que solicita y el solicitado. \\
\textbf{RF:} RF5.\theRF. \textbf{RD(s):} RDE5.\theRDE, RDR5.\theRDR \\
\hspace*{1cm} \textit{Descripción:} Si se desea ver mensajes, deben de existir ambos usuarios, si no, no mostramos nada. \\ [6pt]
\stepcounter{RS}

\stepcounter{RF}
\stepcounter{RDE}
\stepcounter{RDR}
\stepcounter{RDW}
\stepcounter{RDS}
% ======================= %
%      ARCHIVAR CHAT      %
% ======================= %
\section{Des/archivar usuario}
\textbf{\textcolor{rojo}{\underline{RF5.\theRF:}}}. Des/archivar usuario: Permite activar o desactivar la visualización de un usario al listar usuarios. 
Se requiere del ID de ambos usuarios implicados en la conversación (Cadenas de 40 caracteres). Para ello, debe de existir 
el usuario que se desea archivar, y una conversación entre él y el solicitante de archivar el chat. Cuando un chat esté archivado, no se mostrará en la acción ver 
conversaciones. Si se solicita la acción de archivar, se almacena el ID del usuario que archiva y el archivado, \\[6pt]

\textbf{\textcolor{rojo}{Entrada:}} Agente externo: usuario. Acción: solicitar la des/archivación de un usuario. 
Requisito de datos de entrada \textcolor{azul}{RDE5.\theRDE}. \\[6pt]

\textbf{\textcolor{rojo}{BD:}} Requisito de datos de lectura \textcolor{azul}{RDR5.\theRDR} y de escritura \textcolor{azul}{RDW5.\theRDW}. \\[6pt]

\textbf{\textcolor{rojo}{Salida:}} Agente externo: usuario. Acción: confirmación resultado. 
Requisito de datos de salida: ninguno \\[12pt]

\textcolor{azul}{\textbf{RDE5.\theRDE:}} Usuario que deseamos des/archivar   \\
\hspace*{1cm} ID Usuario: Cadena de caracteres (40) \\[12pt]

\textcolor{azul}{\textbf{RDR5.\theRDR:}} Comprobación de que existe el usuario a archivar y un mensaje entre ambos usuarios, además de si el usuario está ya archivado\\
\hspace*{1cm} ID Usuario a archivar: Cadena de caracteres (40) \\
\hspace*{1cm} ID Usuario que archiva: Cadena de caracteres (40) \\[12pt]

\textcolor{azul}{\textbf{RDW5.\theRDW:}} Datos del usuario que se desea archivar/desarchivar, si el ya está archivado, se borra la tupla, si no, se añade. \\
\hspace*{1cm} ID Usuario a archivar: Cadena de caracteres (40) \\
\hspace*{1cm} ID Usuario que archiva: Cadena de caracteres (40) \\[12pt]

\setcounter{RS}{1}
\textbf{\textcolor{verde}{Restricciones Semánticas:}} \\[6pt]
\textbf{\textcolor{verde}{\underline{RS5.\theRF.\theRS:}}} Para archivar un usuario, debe de existir. \\
\textbf{RF:} RF5.\theRF. \textbf{RD(s):} RDE5.\theRDE, RDR5.\theRDR \\
\hspace*{1cm} \textit{Descripción:} Si se desea des/archivar una usuario, el usuario con el que se habla debe de existir, si no, se da un mensaje de error indicando que el usuario no existe y no modificando la base de datos. \\ [6pt]
\stepcounter{RS}

\textbf{\textcolor{verde}{\underline{RS5.\theRF.\theRS:}}} Para archivar un usuario, deben ser amigos recíprocamente, es decir, capaces de tener una conversación. \\
\textbf{RF:} RF5.\theRF. \textbf{RD(s):} RDE5.\theRDE, RDR5.\theRDR \\
\hspace*{1cm} \textit{Descripción:} Si se desea des/archivar un usuario, tienen que ser amigos recíprocamente, de no ser así, se mostrará un error y se indicará que no se puede realizar dicha acción, no modificando la base de datos. \\ [6pt]
\stepcounter{RS}

\stepcounter{RF}
\stepcounter{RDE}
\stepcounter{RDW}
\stepcounter{RDR}

